%قبل از شروع مسئله ترجیحاً چیزی ننویسید. اما اگر خواستید بنویسید، در قالب لیستی (فقرات و شمارش) بنویسید.
%برای اضافه کردن نمره، بعدی از شروع‌مسئله، [x] را اضافه کنید و به جای x  نمره مورد نظر را قرار دهید. برای زیرمسئله‌ها هم به همین صورت است.
%\شروع‌مسئله[10]

\شروع‌مسئله [10+10]
% صورت اصلی مسئله را در این‌جا بنویسید.

%اگر مسئله چند قسمت داشت به این صورت عمل کنید:
\شروع{زیرمسائل}
\زیرمسئله
داده ساختاری طراحی کنید که در زمان $O(1)$ عنصر $i$ ام لیست را بیابد و اعمال درج و حذف عناصر ابتدا و انتها را زمان سرشکن $O(1)$ انجام دهد.

\زیرمسئله
با ادغام داده ساختارهای صف و پشته، داده ساختاری ساختیم که عملیات زیر را پشتیبانی می‌کند.
\شروع{فقرات}
\فقره
:Pop حذف کردن یک عنصر از ابتدای سمت راست لیست. 
\فقره
:Push اضافه کردن یک عنصر به ابتدای سمت راست لیست.
\فقره
:Pull حذف کردن یک عنصر از ابتدای سمت چپ لیست.
\پایان{فقرات}
با استفاده از تنها ۳ پشته و $O(1)$ حافظه اضافی داده ساختار را طوری طراحی کنید که اعمال گفته‌شده در زمان سرشکن $O(1)$ قابل انجام باشد.

\پایان{زیرمسائل}
%اگر مسئله چند قسمت داشت به این صورت عمل کنید:
%\شروع{زیرمسائل}
%\زیرمسئله[0] % مسئلۀ ۱آ
%\زیرمسئله % مسئلۀ ۱ب
%\زیرمسئله % مسئلۀ ۱پ
%\زیرمسئله % مسئلۀ ۱ت


\solution{
\شروع{فقرات}
\فقره
برای این سوال راه حل ها متفاوتی وجود دارد. برای ذخیره سازی $n$ عنصر یک آرایه به طول $3n$ را در نظر می گیریم و عناصر را در وسط آن قرار می‌دهیم. به این ترتیب عملیات پیدا کردن عنصر در زمان $O(1)$ انجام پذیر است همچنین برای درج عنصر در ابتدا کافیست عنصر را به سمت چپ داده‌ها اضافه کنیم. به همین ترتیب می توانیم عناصر را در انتها داده‌ها اضافه کنیم. اگر هنگام درج فضای اضافه‌ای وجود نداشت یک آرایه جدید به طول سه برابر تعداد عناصر کنونی درست می کنیم و آنها را در وسط آرایه جدید قرار می دهیم و آرایه قبلی را حذف می کنیم. توجه کنید تعداد درج‌هایی که در آرایه قبلی انجام داده‌ایم قبل از این‌که آن را بازسازی کنیم خطی بوده است. همچنین برای حذف هنگامی که تعداد عناصر از یک ششم طول آرایه کمتر شد دوباره همین کار را تکرار می‌کنیم. توجه کنید تعداد حذف‌هایی که در آرایه قبلی پیش از بازسازی آن انجام شده است خطی می‌باشد. به این ترتیب اعمال درج و حذف عناصر ابتدا و انتها را زمان سرشکن $O(1)$ قابل انجام است.

ذکر این نکته لازم است که برای دسترسی $O(1)$ به عناصر کافی است اندیس اولین عنصر و تعداد عناصر را نگه‌دار و به روز کنیم.
\فقره
نصف ابتدای عناصر را در یک پشته و نصف انتها را به صورت برعکس در یک پشته دیگر نگه داری می کنیم به صورتی که عناصر اول و آخر در بالای پشته‌ها قرار گیرند. در این صورت عملیات Push در $O(1)$ قابل انجام است. هنگام Pop و Pull اگر هر دو پشته پر باشند می‌توانیم این عملیات را در $O(1)$ انجام دهیم. اگر یکی از پشته‌ها خالی شد (به طور مثال پشته دوم) نصف عنصرهای پشته پر (به طور مثال پشته اول) را Pop کرده و در پشته سوم می ریزیم. سپس از پشته اول دوباره نصف باقی مانده را Pop می‌کنیم و در پشته دوم می ریزیم. سپس تمام عناصر موجود در پشته سوم را Pop می‌کنیم و در پشته اول می ریزیم به این ترتیب عنصر ابتدا و انتها دوباره در بالای پشته دوم و اول قرار می گیرند. همچنین تعداد عملیات انجام شده قبل از هر بازسازی خطی است در نتیجه می توان تمام این عملیات را در زمان سرشکن $O(1)$ انجام داد.
\پایان{فقرات}
}




%%%%%%%%%%%%%%%%%%%%%%%%%%%%%%%%%%%%%%%%%%%%%%%%%%%%%
% برخی دستورات مفید برای تمرین‌ها. %
%%%%%%%%%%%%%%%%%%%%%%%%%%%%%%%%%%%%%%%%%%%%%%%%%%%%%

% فرمول‌های نمونه:
%$f(x) = \Theta(x)$
%$T(x, y) \leq \log(x) + 2^y + \binom{2n}{n}$
% {\tt code\_function}


% لیست موارد بدون شماره‌گذاری:
%\شروع{فقرات}
%    \فقره مورد اول 
%        \شروع{فقرات}
%            \فقره مورد اول داخلی 
%                \شروع{فقرات}
%                    \فقره اولین مورد خیلی داخلی 
%                    \فقره دومین مورد خیلی داخلی 
%                \پایان{فقرات}
%            \فقره مورد دوم داخلی 
%        \پایان{فقرات}
%    \فقره مورد دوم 
%\پایان{فقرات}

% لیست موارد شماره‌گذاری شده:
%\شروع{شمارش}
%    \فقره یک مورد 
%    \فقره یک مورد دیگر 
%    \فقره باز هم یک مورد دیگر
%\پایان{شمارش}

% برای نوشتن یک فرمول چند خطی که یک شماره بخوره. برای حالت‌های دیگر به سایت زیر مراجعه کنید:
% https://www.overleaf.com/learn/latex/Aligning_equations_with_amsmath
%\begin{align} 
%    \begin{split}
%        (x+y)^3 &= (x+y)^2(x+y) \\
%                &= (x^2+2xy+y^2)(x+y) \\
%                &= (x^3+2x^2y+xy^2) + (x^2y+2xy^2+y^3) \\
%                &= x^3+3x^2y+3xy^2+y^3
%    \end{split}                                 
%\end{align}

% نوشتن ماتریس و آرایه دوبعدی
%\begin{align}
%    A = 
%    \begin{bmatrix}
%        A_{11} & A_{21} \\
%        A_{21} & A_{22}
%    \end{bmatrix}
%\end{align}

% اضافه کردن تصویر با شرح. برای اضافه کردن تصویر بدون شرح، فقط دستور درج‌تصویر را استفاده کنید.
%\شروع{شکل}[h!]
%\تنظیم‌ازوسط
%  \درج‌تصویر[پهنا=0.5\پهنای‌متن]{img.jpg}
%\شرح{شرح تصویر}
%\برچسب{تصویر:۱}
%\پایان{شکل}
%\پایان{زیرمسائل}

%کد زیر را ببینید.
%%\lstinputlisting{source_filename.py}
%\begin{latin}
%\begin{lstlisting}
%# Import the modules
%import sys
%import random
%
%ans = True
%
%while ans:
%question = raw_input("Ask the magic 8 ball a question: (press enter to quit) ")
%
%answers = random.randint(1,8)
%
%if question == "":
%sys.exit()
%
%elif answers == 1:
%print "It is certain"
%
%elif answers == 2:
%print "Outlook good"
%
%elif answers == 3:
%print "You may rely on it"
%
%elif answers == 4:
%print "Ask again later"
%
%elif answers == 5:
%print "Concentrate and ask again"
%
%elif answers == 6:
%print "Reply hazy, try again"
%
%elif answers == 7:
%print "My reply is no"
%
%elif answers == 8:
%print "My sources say no"
%\end{lstlisting}
%\end{latin}
